\documentclass[11pt]{article}
\usepackage[a4paper,margin=1.9cm]{geometry}
\usepackage[T1]{fontenc}
\usepackage{textcomp}
\usepackage[spanish]{babel}
\usepackage{amsmath,amssymb}
\usepackage{lmodern}
\renewcommand{\familydefault}{\sfdefault}
\usepackage{enumitem}

\newcommand{\vect}[2]{\begin{pmatrix}#1\\#2\end{pmatrix}}
\usepackage{tikz}
\usetikzlibrary{calc,arrows.meta,patterns,decorations.pathmorphing}
\usepackage[table]{xcolor}
\usepackage{multicol}
\usepackage{parskip}

\definecolor{unalblue}{RGB}{31,73,124}

\setlist[enumerate,1]{itemsep=1.4em,topsep=1em}

\newcommand{\examheader}{%
  \noindent
  \begin{minipage}[t]{0.6\linewidth}
    {\small\bfseries Geometría Vectorial y Analítica --- Parcial 1}\\
    {\small Semestre 2015--01}
  \end{minipage}%
  \hfill
  \begin{minipage}[t]{0.35\linewidth}
    \raggedleft\small Escuela de Matemáticas. Universidad Nacional de Colombia.
  \end{minipage}\\[2pt]
  \rule{\linewidth}{0.6pt}
}
\newcommand{\examfooter}{%
  \vfill
  \rule{\linewidth}{0.4pt}\\[1pt]
  {\scriptsize\textit{Transcripción digital --- MathUNAL}\hfill\textcolor{unalblue}{\textbf{\thepage}}}
}
\pagestyle{empty}

\newcommand{\nota}[1]{{\small\itshape\color{unalblue!70!black} #1}}

\begin{document}
\examheader

\begin{center}
{\Large Primer Examen Parcial de Geometría Vectorial}\\[4pt]
{\small 7 de marzo de 2015}
\end{center}

\vspace{10pt}
\textbf{PUNTAJE. SOLO PARA USO OFICIAL}

\begin{center}
\begin{tabular}{|c|c|c|c|c|}
\hline
1 & 2 & 3 & 4 & TOTAL \\
\hline
 & & & & \\
\hline
\end{tabular}
\end{center}

\vspace{10pt}
\textbf{INSTRUCCIONES}

No está permitido el uso de calculadora ni de otros aparatos electrónicos. Solucione las preguntas en los espacios en blanco asignados a cada una de ellas. No está permitido sacar hojas en blanco ni ningún tipo de apuntes durante el examen. Utilice la página 6 para realizar sus operaciones en borrador. \textbf{Duración: 1 hora y 50 minutos.}

\begin{enumerate}
\item Considere la recta $\mathcal{L}: 2x-y-3=0$ y el punto $U=\vect{2}{3}$.

\begin{enumerate}[label=(\alph*)]
\item (10\%) Halle la proyección ortogonal de $U$ sobre $\mathcal{L}$, $\mathrm{P}_{\mathcal{L}}(U)$.
\item (15\%) Encuentre la ecuación de la circunferencia que tiene centro en $U$ y pasa por $\mathrm{P}_{\mathcal{L}}(U)$.
\end{enumerate}

\item Considere el punto $U=\vect{3}{-2}$ y la recta $\mathcal{L}$ con ecuaciones paramétricas
$$\mathcal{L}:\begin{cases}x=-2+3t,\\ y=5+t,\end{cases}\quad t\in\mathbb{R}.$$

\begin{enumerate}[label=(\alph*)]
\item (5\%) Pruebe que $U$ no pertenece a $\mathcal{L}$.
\item (10\%) Encuentre una ecuación general de la recta $\mathcal{L}_1$ que es paralela a $\mathcal{L}$ y pasa por $U$.
\item (10\%) Encuentre una ecuación general de la recta $\mathcal{L}_2$ que es perpendicular a $\mathcal{L}$ y pasa por $U$.
\item (10\%) Halle el punto $Q$ de la recta $\mathcal{L}$ que está más cercano al punto $U$.
\end{enumerate}

\item (25\%) Sea $\mathcal{L}_1$ la recta que pasa por $P=\vect{-1}{-1}$ y es perpendicular a la recta $\mathcal{L}_2: x-2y-5=0$. Encuentre el punto de intersección de $\mathcal{L}_1$ y la recta $\mathcal{L}_3: 7x-3y-25=0$.

\item (15\%) Sean $X$ y $U$ vectores del plano. Usando la definición y las propiedades del producto punto entre vectores, pruebe que
$$\|X+U\|^2+\|X-U\|^2=2\|X\|^2+2\|U\|^2.$$

\end{enumerate}

\examfooter
\end{document}
